\lecture{5}{Fri 12 Sep 2025 12:21}{Tangent vectors}

Define $\mathbb{R} ^n_a$ as the set $\left\{ a \right\} \times \mathbb{R} ^n$ for some vector $a\in\mathbb{R} ^n$. We will show that this admits an interpretation as a tangent space to $\mathbb{R} ^n$ at $a$. We say that an element $(a,v)\in\mathbb{R} _a^n$ is a \emph{geometric tangent vector}, and we write it as $v_a$ or $v|_a$. The intuitive interpretation of this object is taking the vector $v$, and giving it starting point $a$.

The set $\mathbb{R} _a^n$ forms a vector space in the obvious way, by interpreting $\left\{ a \right\} $ as $0$, and simply working with the space $0\times \mathbb{R} ^n\cong \mathbb{R} ^n$. The vectors $e_i|_a$ thus form a basis for this space.

This allows us to define tangent spaces to manifolds which admit embeddings into $\mathbb{R}^{n+1} $, by taking the tangent space at a point $a\in M$ as $\mathbb{R} ^n_a$. However, to generalize this definition arbitrarily, we must explore some alternative representations of this information.

To do this, we must somehow translate this information into that of smooth maps and manifolds. For example, given any vector $v_a\in\mathbb{R} _a^n$, the directional derivative $D_v|_a : C^\infty (\mathbb{R} ^n) \to \mathbb{R} $ given by
\[
	D_v|_af=D_vf(a)=\left.\frac{\mathrm{d}}{\mathrm{d}t} \right|_{t=0} f(a+tv)
\]
is $\mathbb{R} $-linear, and satisfies the \emph{product rule}
\[
	D_v|_a(fg)=f(a)D_v|_ag+g(a)D_v|_af
.\]
If we write $v_a=v^ie_i|_a$, then $D_v|_af$ can be written more concretely as
\[
	D_v|_af=v^i \frac{\partial f}{\partial x^i} (a)
.\]
This follows by the chain rule. We generalize directional derivatives using this formulation.

\begin{definition}\label{1.2.7}
	Let $a\in\mathbb{R} ^n$. A \emph{derivation} at $a$ is an $\mathbb{R} $-linear map $w : C^\infty (\mathbb{R} ^n) \to \mathbb{R} $ satisfying the product rule
	\[
		w(fg)=f(a)w(g)+g(a)w(f)
	.\]
	We write $T_a\mathbb{R} ^n$ for the set of all derivations at $a$.
\end{definition}

We immediately see that $T_a\mathbb{R} ^n$ is a vector space by noting that if $w_1,w_2$ are derivations, then
\[
	(w_1+w_2)(fg)=(f(a)w_1(g)+g(a)w_1(f))+(f(a)w_2(g)+g(a)w_2(f))=f(a)(w_1+w_2)(g)+g(a)(w_1+w_2)(f)
,\]
\[
	(cw)(fg)=c\left( f(a)w(g)+g(a)w(f) \right) =f(a)(cw)(g)+g(a)(cw)(f)
.\]
The following lemma is not too interesting, but implies a very interesting result.

\begin{lemma}\label{1.2.8}
	Let $a\in\mathbb{R} ^n$, $w\in T_a\mathbb{R} ^n$, and $f,g\in C^\infty (\mathbb{R} ^n)$. If $f$ is a constant function, then $w(f)=0$. Additionally, if $f(a)=g(a)=0$, then $w(fg)=0$.
\end{lemma}
\begin{proof}
	For the first statement, it suffices to prove $f$ for the constant function $f(x)=1$, since $g(x)=c$ implies $g=cf$, and the statement follows by linearity of derivations. The product rule gives
	\[
		w(f)=w(ff)=1w(f)+1w(f)=2w(f)
	\]
	and thus $w(f)=0$.

	For the second statement, we have
	\[
		w(fg)=f(a)w(g)+g(a)w(f)=0w(f)+0w(g)=0
	.\]
\end{proof}

\begin{proposition}\label{1.2.9}
	Let $a\in\mathbb{R} ^n$ and $v_a\in\mathbb{R} _a^n$. Then $D_v|_a$ is a derivation, and the map $v_a \mapsto D_v|_a$ is an isomorphism of $\mathbb{R} $-vector spaces $\mathbb{R} _a^n \cong T_a\mathbb{R} ^n$.
\end{proposition}
\begin{proof}
	First, we show $D_v|_a$ is a derivation. This follows by the properties of the partial derivative $\partial _\mu $. To show it is $\mathbb{R} $-linear, we have
	\[
		D_v|_a(f+g)=v^\mu \partial _\mu (f+g)|_a=v^\mu \partial _\mu f|_a+v^\mu \partial _\mu g|_a=D_v|_af+D_v|_ag,
	\]
	\[
		D_v(cf)=v^\mu \partial _\mu (cf)|_a=cv^\mu \partial _\mu f|_a=cD_v|_af
	.\]
	The product rule also follows by the fact that $\partial _\mu $ satisfies the product rule: $\partial _x(fg)+f \partial _x(g)+g \partial _xf$. Adding in the evaluation gives
	\[
	D_v|_a(fg)=v^\mu \left( f(a)\partial _\mu g|_a+g(a)\partial _\mu f|_a \right) =f(a)D_v|_ag+g(a)D_v|_af
	.\]
	
	Now we must show the map $v_a\mapsto D_v|_a$ is $\mathbb{R} $-linear. To show additivity, we must show $D_{v+w} |_a=D_v|_a+D_w|_a$. This is obvious by the fact that $(v+w)^\mu =v^\mu +w^\mu $:
	\[
		D_{v+w} |_a=(v+w)^\mu \partial _\mu |_a=v^\mu \partial _\mu |_a+w^\mu \partial _\mu |_a
	.\]
	Similarly, since $(cv)^\mu =cv^\mu $, we have
	\[
		D_{cv} |_a=(cv)^\mu \partial _\mu |_a=cv^\mu \partial _\mu |_a=cD_v|_a
	.\]
	Therefore, the map is $\mathbb{R} $-linear.

	The main part of the proof is that the map is an isomorphism of $\mathbb{R} $-vector spaces. Injectivity is show by showing the kernel is trivial. Since we do not know that $T_a\mathbb{R} ^n$ is finite dimensional, this does not suffice as an isomorphism proof. Let $D_v|_a$ be the zero derivation $f\mapsto 0$. We want to show $v_a=0$. Let $f$ be the coordinate function $x^j : \mathbb{R} ^n \to \mathbb{R} $ by projecting onto the $j$th component. Then $\partial _ix^j=\delta_i^j$, giving
	\[
		0 = D_v|_a = v^\mu \partial _\mu x^j|_a=v^j ,
	\]
	since $1|_a=1$. The statement follows by letting $j$ range over $[n]$, showing the map is injective.

	To prove surjectivity, let $w\in T_a\mathbb{R} ^n$ be an arbitrary derivation. With the same definition of $x^j$, write $v^i=w(x^i)$. We show $w=D_v|_a$.

	We show this by proving they act on maps $f : C^\infty (\mathbb{R} ^n)$ in the same way. Since $f$ is analytic, it admits a Taylor expansion
	\[
		f=f(a)+\sum_{i=1}^{n} \frac{\partial f}{\partial x^i} (a)\left( x^i-a^i \right) +\sum_{i,j=1}^{n} \left( x^i-a^i \right) \left( x^j-a^j \right) \int_{0}^{1} (1-t)\frac{\partial ^2f}{\partial x^i \partial x^j} (a+t(x-a)) \,\mathrm{d} t
	.\]
	The functions $x^i-a^i$ and $x^j-a^j$ vanish at $x=a$. Since $w$ is $\mathbb{R} $-linear and distributes over addition, by \cref{1.2.8},
	\[
		w((x^i-a^i)(x^j-a^j))=0
	.\]
	The derivation thus annihilates the last term:
	\[
		w(f)=w(f)(a)+\sum_{i=1}^{n} w\left( \frac{\partial f}{\partial x^i} (a)\left( x^i-a^i \right)  \right) =\sum_{i=1}^{n} \frac{\partial f}{\partial x^i} (a)\left( w(x^i)-w(a^i) \right) =\sum_{i=1}^{n} \frac{\partial f}{\partial x^i} (a)w(x^i)
	.\]
	Since $w(x^i)=v^i$, this is simply $v^\mu \partial _\mu f|_a=D_v|_a$. This proves surjectivity.
\end{proof}

\begin{corollary}\label{1.2.10}
	$T_a\mathbb{R} ^n$ has dimension $n$, and has a basis given by
	\[
		\left.\frac{\partial }{\partial x^i}\right|_a \coloneqq D_{e_i} |_a
	\]
	In particular, these are defined by
	\[
		\left.\frac{\partial }{\partial x^i} \right|_af=\frac{\partial f}{\partial x^i} (a)
	.\]
\end{corollary}
\begin{proof}
	These follow immediately by the above proof.
\end{proof}

We now see an obvious way to generalize this definition to arbitrary manifolds. We can define derivations on arbitary $C^\infty $-manifolds, and then show that this is generated by the partials at a point.

\begin{definition}\label{1.2.11}
	Let $M$ be a smooth manifold, and let $p\in M$. A linear map $v : C^\infty (M) \to \mathbb{R} $ is a \emph{derivation at p} if it satisfies the product rule
	\[
		\forall f,g\in C^\infty (\mathbb{R} ) \, : \, v(fg)=f(p)v(g)=g(p)v(f)
	.\]
	The set of all derivations for $p$ is denoted by $T_pM$, and is called the \emph{tangent space} to $M$ at $p$. An element of $T_pM$ is a \emph{tangent vector} at $p$.
\end{definition}

We now prove some facts to show that this definition behaves as expected.

\begin{lemma}\label{1.2.12}
	Let $M$ be a smooth manifold with or without boundary, $p\in M$, $v\in T_pM$, and $f,g\in C^\infty (M)$. If $f$ is constant, then $v(f)=0$. If $f(p)=g(p)=0$, then $v(fg)=0$.
\end{lemma}
\begin{proof}
	The same proof for \cref{1.2.8} suffices.
\end{proof}

